\subsection{Selection and Documentation of Measurement Method}

After the quantity to be measured has been established and the appropriate instruments selected, the experimenter should prepare documentation of the measurement principle. In metrology, this is referred to as a **measurement procedure**, or simply a procedure. This is one of the most important tasks of a measurement engineer: documenting the work in sufficient detail so that it can be reproduced with the same results and the same order-of-magnitude uncertainties.

Preparing a measurement procedure is a rigorous task and requires specifying:

\begin{enumerate}
\item What is being measured.
\item Which instruments are used.
\item How stable the system is (e.g., noise, drift, environmental effects).
\item The estimated measurement uncertainty.
\item How everything is connected (a diagram or picture guide is highly recommended).
\end{enumerate}

This is non-negotiable: if data collection is to be performed, the engineer \textbf{must} present a measurement procedure to both the supervisor and laboratory technicians. This ensures that the correct method is used and prevents damage to instruments or personnel.

It should be noted that in experimental work, the requirements for uncertainty estimation may not be as stringent as in industry. However, the experimenter should always remain aware of potential sources of uncertainty and make rough estimates to ensure that the measured quantity is not overwhelmed or obscured by unknown factors.

The procedure should also describe how data is recorded and documented, whether it is collected manually or through an automated system. If automation is used, the code or scripts responsible for data acquisition should be stored alongside the procedure to ensure reproducibility and traceability.


A separate guide will provide more detailed instructions on designing an appropriate measurement procedure. This current section focuses on the principles and documentation requirements, while the other guide offers practical templates and step-by-step guidance.

It should be noted that the measurement procedure can apply to any system, even something as simple as collecting data using an Arduino and its ADC. In such cases, the procedure should describe how the device is used, how it has been verified or calibrated, and any relevant settings or connections. The purpose of the procedure is to ensure that the measurements can always be reproduced and, if necessary, traced back for verification or analysis.
